problem:  
Let \(M^6\) be a **compact simply connected** smooth 6‑manifold that admits both an *integrable* SU(3)‑structure (Calabi–Yau) \((\omega_0, \psi_{+,0}, \psi_{-,0})\) and a *strict nearly Kähler* SU(3)‑structure \((\omega_1, \psi_{+,1}, \psi_{-,1})\), possibly inducing different metrics.  
Such coexistence is not known to occur on any compact 6‑manifold, and the two types of structures impose radically different curvature and topological conditions: the Calabi–Yau is Ricci‑flat with vanishing first Chern class, while the nearly Kähler is Einstein with positive scalar curvature and non‑zero torsion class \(W_1\).  

**Problem:**  
Does there exist a compact simply connected 6‑manifold \(M\) that supports *both* a strict nearly Kähler SU(3)-structure and an integrable (Calabi–Yau) SU(3)-structure?  
If coexistence is either possible or impossible, determine which of the following mechanisms governs it:

1. **Topological obstruction:** Does the coexistence of NK and CY structures contradict characteristic class identities, e.g. the incompatibility between \(p_1(M)\), \(c_1(M)\), or decomposability properties of \(H^{2,0}(M)\)?  
2. **Geometric–analytic obstruction:** Is the Einstein condition and positive scalar curvature of the NK metric analytically incompatible with Ricci‑flat structures on the same smooth manifold (for instance via Hitchin functionals or volume‑normalization constraints)?  
3. **Hypothetical coexistence:** If such an \(M\) could exist (even in a nonhomogeneous setting), could the two SU(3)‑structures be related by deformation, rescaling, or singular limit phenomena, such as a degeneration of the intrinsic torsion to zero?  

---

Why is it a "good" problem:  
This form isolates the truly *fundamental* unknown: whether the same compact manifold can realize both ends of the SU(3) holonomy spectrum—torsion‑free (Calabi–Yau) and strictly torsionful (nearly Kähler). It avoids speculative path questions and instead asks for the geometric, topological, or analytic *obstructions* to coexistence.  

The question bridges topology (characteristic classes and cohomology restrictions), Riemannian holonomy theory, and nonlinear PDE geometry (Einstein and torsion‑free equations). Its resolution would either identify a deep incompatibility theorem for special holonomy types on the same manifold or uncover completely new examples connecting Calabi–Yau and nearly Kähler geometries.  

It is simple to state—“Can one compact 6‑manifold carry both Calabi–Yau and strict nearly Kähler SU(3)-structures?”—yet its implications would be profound for the classification of SU(3)-structures, making it a genuinely *good* and fundamental research‑level problem.